\documentclass[10pt,a4paper]{article}
\usepackage[T1]{fontenc}
\usepackage[margin=2cm]{geometry}
\usepackage{isabelle,isabellesym}
\usepackage{amsmath,amssymb}
\usepackage{pdfsetup}

\urlstyle{rm}
\isabellestyle{it}

\begin{document}

\title{Consecutive Coprime Divisors}
\author{Carles Mar\'in}
\maketitle

\begin{abstract}
For a natural number $n$ let $\tau_\perp(n)$ be the number of pairs of consecutive divisors of $n$, in increasing order, that are
coprime. Erd\H{o}s and Hall~\cite{ErdosHall1978} note that $\tau_\perp(n)=\omega(n)$ when $n=p_1\cdots p_k$ with
$p_i>p_1\cdots p_{i-1}$ for $2\le i\le k$. We formalize this: for every $k$ there is a squarefree $n$ with exactly $k$ distinct prime
factors and $\tau_\perp(n)=k$. The proof describes the increasing list of divisors of such an $n$ by induction on the primes.
\end{abstract}

\paragraph{Acknowledgement.} The formalization and the text of this document were prepared with the
assistance of generative AI systems, and checked by the author.

\tableofcontents

\parindent 0pt\parskip 0.5ex

\input{session}

\bibliographystyle{abbrv}
\bibliography{root}

\end{document}
